%! Function Model Selection and Construction — Unit 1, Lesson 1.7 — Exit Ticket (Answer Key)
\documentclass[10pt]{article}
\usepackage{precalculus-article}
\usepackage{precalculus-key}   % pulls in -boxes; do NOT also load -boxes
\newcommand{\TallMath}[1]{$\displaystyle #1\rule[-1.4em]{0pt}{3.2em}$}

\begin{document}

\pageheader{Unit 1, Lesson 1.7}{Exit Ticket --- Answer Key}

\vspace{0.12in}

\begin{enumerate}[itemsep=10pt]

\item The question is the \ans{$x$} row. Take differences until they are
      constant, then name the model type.

\smallskip
\renewcommand{\arraystretch}{1.4}
\begin{center}
\begin{tabular}{|c|c|c|c|c|c|}
\hline
$x$ & 0 & 1 & 2 & 3 & 4 \\
\hline
$y$ & 2 & 5 & 14 & 29 & 50 \\
\hline
1st difference & --- & \ans{3} & \ans{9} & \ans{15} & \ans{21} \\
\hline
2nd difference & --- & --- & \ans{6} & \ans{6} & \ans{6} \\
\hline
\end{tabular}
\end{center}

\smallskip
Model type: \ans{quadratic}

\item A dog run is built against the side of a garage using \textbf{40 feet} of
      fencing on the other three sides. Take the \textbf{width} $w$ as your
      question (the two fenced ends); $\ell$ is the length along the garage.

\begin{work}
   2w + \ell &= 40 \\
        \ell &= 40 - 2w \\
        A(w) &= w(40 - 2w) \\
             &= 40w - 2w^2
\end{work}

\begin{enumerate}[label=(\alph*), itemsep=4pt]
  \item Constraint: \ans{2} $w \; + \;$ \ans{$\ell$} $= 40$
  \item $A(w) = $ \ans{$40w - 2w^2$}
  \item Domain: \ans{0} $< w <$ \ans{20}
  \item Write one assumption this model makes about the real world.

        \ansline{All 40 ft is used, the run is exactly rectangular, and the garage is long enough.}
\end{enumerate}

\item Two rectangular rugs. In both, the question is the width $w$ and the
      answer is the area.

\begin{enumerate}[label=(\alph*), itemsep=4pt]
  \item The length is always \textbf{3 times} the width: \quad
        $A(w) = $ \ans{$3w^2$}, \quad model type \ans{quadratic}
  \item The length is \textbf{fixed at 8 feet}: \quad
        $A(w) = $ \ans{$8w$}, \quad model type \ans{linear}
  \item Both answers are areas. Why is only one of them quadratic?

        \ansline{In (a) both lengths grow with the width; in (b) only one does.}
\end{enumerate}

\end{enumerate}

\end{document}
